\chapter{Conclusão}

Este capítulo apresenta as conclusões do trabalho os as considerações finais.
Elas devem ser apresentadas de forma clara, concisa e lógica com base nos resultados obtidos ao longo do trabalho. 

É importante que este capítulo converse bem com o capítulo de introdução.
Ou seja, as conclusões devem ser fundamentadas nos objetivos formulados na introdução e nas comprovações discutidas ao longo do trabalho.
Uma sugestão é que o autor ao finalizar o trabalho leia o capítulo de introdução seguido do capítulo de conclusão (ou considerações finais) e analise se está mantendo uma sequência lógica de raciocínio.

A conclusão mantém o foco em fechar o argumento principal do trabalho, sintetizando as descobertas. 
Assim, retoma os objetivos, apresenta resultados-chave e responde à problemática.
Possui natureza de indicar uma verdade mais final ou definitiva sobre a pesquisa.
Já as considerações finais representam uma reflexão mais ampla sobre as descobertas com foco em refletir criticamente sobre as ideias desenvolvidas, sem apresentar uma conclusão definitiva.
Ela apresenta ideias com novos pontos de vista, sugere direções futuras e reconhece limitações do estudo.

Um TCC, dependendo do assunto, pode ter apenas conclusões, apenas considerações finais ou ambos.
No caso de ambos, sugere-se que o capítulo tenha o título ``CONCLUSÃO'' e, dentro do capítulo, tenha uma seção de considerações finais.

\section{Sugestões para futuros trabalhos}

Como parte final do TCC, sugere-se incluir recomendações para trabalhos futuros, indicando uma forma de continuação dos estudos envolvendo o tema do trabalho.